\begin{abstract}
We demonstrate that lead-lag conclusions in iterative model collapse are determined more by method choice than by the data: naive cross-correlation yields 100\% significance (16/16 metric pairs), while Toda--Yamamoto Granger causality retains only 25\% (4/16), with inter-method agreement of $\kappa = 0.185$. We identify three methodological pitfalls---method sensitivity, spurious regression under shared trends, and regime dependence---that together inflate spurious-signal rates to 75\%. To replace single-method reporting, we introduce a multi-method consensus framework requiring agreement across ${\geq}4$ of 5 statistical tests with FDR correction, achieving 0\% false-positive and 0\% false-negative rates on reference extreme regimes. Applied to Pythia-410M across 4 seeds, 5 synthetic-data mixing ratios ($\alpha$), and 10 iterative training generations on WikiText-103, the framework reveals a sharp regime boundary: significant lead-lag relationships emerge at $\alpha{=}1.0$ (with preliminary single-seed evidence at $\alpha{=}0.75$), where token entropy Granger-causes diversity collapse with a one-generation lead ($|r| > 0.99$, 24/24 pairs significant, 0/24 in the reverse direction). Differencing eliminates all perplexity leads ($8/8 \to 0/8$) but preserves all diversity leads (24/24), confirming token entropy as a robust early-warning canary metric.
\end{abstract}